\documentclass[11pt,a4paper]{article}
\usepackage[utf8]{inputenc}
\usepackage[T1]{fontenc}
\usepackage[brazilian]{babel}
\usepackage{lmodern}
\usepackage{amsmath,amssymb}
\usepackage[margin=2.25cm]{geometry}
\usepackage{booktabs,array}
\usepackage{xcolor}
\usepackage{fancyhdr}
\usepackage{listings}
\usepackage[hidelinks]{hyperref}

\definecolor{azul}{HTML}{183B56}
\definecolor{cinza}{RGB}{245,245,245}
\definecolor{cinzaEscuro}{RGB}{80,80,80}

\pagestyle{fancy}
\fancyhf{}
\fancyhead[L]{\small Projeto e Análise de Algoritmos -- Prof. Felipe A. Louza}
\fancyhead[R]{\small Lista 7}
\fancyfoot[L]{\small Arthur Pires Lopes -- 12221EEL001}
\fancyfoot[R]{\thepage}
\setlength{\headheight}{14pt}
\setlength{\parindent}{0pt}
\setlength{\parskip}{6pt}
\setlength{\emergencystretch}{2em}
\renewcommand{\arraystretch}{1.3}

\newcommand{\res}[1]{\[\boxed{#1}\]}
\newcommand{\topico}[1]{\par\medskip\textbf{\color{azul}#1}\par}
\newcommand{\questao}[2]{\section*{Questão #1 -- #2}}

\lstset{
  basicstyle=\ttfamily\small,
  backgroundcolor=\color{cinza},
  frame=single,
  rulecolor=\color{black!20},
  columns=fullflexible,
  keepspaces=true,
  showstringspaces=false,
  numbers=left,
  numberstyle=\tiny\color{cinzaEscuro},
  numbersep=7pt,
  literate={á}{{\'a}}1 {ã}{{\~a}}1 {ç}{{\c c}}1 {é}{{\'e}}1
           {í}{{\'i}}1 {ó}{{\'o}}1 {ú}{{\'u}}1 {Á}{{\'A}}1
           {É}{{\'E}}1 {Í}{{\'I}}1 {Ó}{{\'O}}1 {Ú}{{\'U}}1
           {à}{{\`a}}1 {õ}{{\~o}}1 {â}{{\^a}}1 {ê}{{\^e}}1
           {ô}{{\^o}}1 {<=}{{$\le$}}1 {>=}{{$\ge$}}1 {!=}{{$\neq$}}1
           {<-}{{$\leftarrow$}}1
}

\begin{document}

\begin{center}
  {\small Universidade Federal de Uberlândia -- UFU\\Faculdade de Engenharia Elétrica}\par
  \vspace{4pt}
  {\LARGE\bfseries\color{azul} Lista 7: O Problema da Seleção}\par
  \vspace{4pt}
  \textbf{Arthur Pires Lopes}\quad Matrícula: \textbf{12221EEL001}\par
  \textit{Resolução formal para entrega -- Projeto e Análise de Algoritmos (Prof. Felipe A. Louza)}
\end{center}
\hrule
\vspace{6pt}

% ==============================================================================
% QUESTÃO 1
% ==============================================================================
\questao{1}{Mínimo e Máximo Simultâneos com no Máximo $3\lfloor n/2 \rfloor$ Comparações}

\textbf{Enunciado:} Escreva o pseudocódigo de algoritmo que encontre o máximo e o mínimo em um vetor $A[1, n]$ realizando no máximo $3\lfloor n/2 \rfloor$ comparações.

\topico{Estratégia Algorítmica}
Em vez de comparar cada elemento independentemente com o mínimo e o máximo correntes (o que demandaria $2n - 2$ comparações), processamos os elementos \textbf{aos pares}. 
\begin{enumerate}
    \item Para cada par $(A[i], A[i+1])$, realizamos 1 comparação entre seus dois elementos para determinar qual é o menor e qual é o maior do par.
    \item Comparamos o menor do par com o mínimo corrente $\min$ (1 comparação).
    \item Comparamos o maior do par com o máximo corrente $\max$ (1 comparação).
\end{enumerate}
Dessa forma, processamos 2 elementos utilizando exatamente $3$ comparações.

\topico{Pseudocódigo}
\begin{lstlisting}
MIN-MAX(A, n)
 1  se n == 1 então
 2      retorne (A[1], A[1])
 3  se n é ímpar então
 4      min <- A[1]
 5      max <- A[1]
 6      inicio <- 2
 7  senão
 8      se A[1] <= A[2] então
 9          min <- A[1]; max <- A[2]
10      senão
11          min <- A[2]; max <- A[1]
12      inicio <- 3
13  para i <- inicio até n com passo 2 faça
14      se A[i] <= A[i+1] então
15          se A[i] < min então min <- A[i]
16          se A[i+1] > max então max <- A[i+1]
17      senão
18          se A[i+1] < min então min <- A[i+1]
19          se A[i] > max então max <- A[i]
20  retorne (min, max)
\end{lstlisting}

\topico{Contagem Exata de Comparações de Elementos}
\begin{itemize}
    \item \textbf{Se $n$ é ímpar:} Inicializamos $\min = \max = A[1]$ sem nenhuma comparação entre chaves. Restam $n-1$ elementos, formando $\frac{n-1}{2} = \lfloor n/2 \rfloor$ pares. Cada par realiza $3$ comparações no laço:
    \[
    C(n) = 0 + 3 \cdot \left(\frac{n-1}{2}\right) = 3\left\lfloor \frac{n}{2} \right\rfloor.
    \]
    \item \textbf{Se $n$ é par:} Inicializamos comparando $A[1]$ e $A[2]$ ($1$ comparação). Restam $n-2$ elementos, agrupados em $\frac{n-2}{2} = \frac{n}{2} - 1 = \lfloor n/2 \rfloor - 1$ pares. Cada par realiza $3$ comparações:
    \[
    C(n) = 1 + 3 \cdot \left(\frac{n}{2} - 1\right) = \frac{3n}{2} - 2 = 3\left\lfloor \frac{n}{2} \right\rfloor - 2.
    \]
\end{itemize}
Em ambos os casos, o número total de comparações satisfaz:
\res{C(n) \le 3\left\lfloor \frac{n}{2} \right\rfloor.}

\vspace{12pt}
\hrule
\vspace{6pt}
% ==============================================================================
% QUESTÃO 2
% ==============================================================================
\questao{2}{Execução Passo a Passo do Algoritmo \texttt{SELECTION}}

\textbf{Enunciado:} Mostre a execução passo-a-passo do algoritmo \texttt{SELECTION} para encontrar o $3^\circ$ menor elemento para o seguinte vetor:
\[
A = \langle 1, 8, 7, 6, 5, 3, 2, 10, 9 \rangle
\]
Considere que o pivot é escolhido sempre como o primeiro elemento do vetor.

\topico{Definição do Algoritmo e da Partição}
O algoritmo busca a $i$-ésima menor estatística de ordem no subvetor $A[p \dots r]$.
\begin{lstlisting}
SELECTION(A, p, r, i)
 1  se p == r então retorne A[p]
 2  q <- PARTITION-PRIMEIRO(A, p, r)
 3  k <- q - p + 1   // rank relativo do pivô no subvetor
 4  se i == k então retorne A[q]
 5  senão se i < k então retorne SELECTION(A, p, q - 1, i)
 6  senão retorne SELECTION(A, q + 1, r, i - k)
\end{lstlisting}
A sub-rotina \texttt{PARTITION-PRIMEIRO(A, p, r)} fixa o pivô $x = A[p]$, mantém o índice de fronteira $m = p$ para elementos menores ou iguais a $x$, percorre $j$ de $p+1$ até $r$, trocando $A[++m]$ com $A[j]$ sempre que $A[j] \le x$, e por fim troca $A[p]$ com $A[m]$, retornando a posição final $q = m$.

\topico{Rastreamento Passo a Passo}
Buscamos o $3^\circ$ menor elemento no vetor global com $n = 9$, ou seja, chamada inicial \texttt{SELECTION(A, 1, 9, 3)}.

\begin{enumerate}
    \item \textbf{Chamada 1:} \texttt{SELECTION(A, 1, 9, 3)}
    \begin{itemize}
        \item Subvetor: $A[1 \dots 9] = \langle \mathbf{1}, 8, 7, 6, 5, 3, 2, 10, 9 \rangle$, buscando $i = 3$.
        \item Pivô: $x = A[1] = 1$.
        \item Varredura ($j = 2 \dots 9$): todos os elementos $\{8, 7, 6, 5, 3, 2, 10, 9\}$ são maiores que $1$. Nenhuma troca ocorre; $m = 1$.
        \item Troca do pivô: $A[1] \leftrightarrow A[1]$. Posição final do pivô: $q = 1$.
        \item Estado do vetor: $\langle \underline{1}, 8, 7, 6, 5, 3, 2, 10, 9 \rangle$.
        \item Posto relativo: $k = q - p + 1 = 1 - 1 + 1 = 1$.
        \item Como $i = 3 > k = 1$: o elemento procurado está à direita. Nova ordem: $i' = 3 - 1 = 2$.
    \end{itemize}

    \item \textbf{Chamada 2:} \texttt{SELECTION(A, 2, 9, 2)}
    \begin{itemize}
        \item Subvetor: $A[2 \dots 9] = \langle \mathbf{8}, 7, 6, 5, 3, 2, 10, 9 \rangle$, buscando $i = 2$.
        \item Pivô: $x = A[2] = 8$. Índice $m$ inicializa em $2$.
        \item Varredura ($j = 3 \dots 9$):
        \begin{itemize}
            \item $j=3: A[3]=7 \le 8 \implies m=3$, troca $A[3] \leftrightarrow A[3]$.
            \item $j=4: A[4]=6 \le 8 \implies m=4$, troca $A[4] \leftrightarrow A[4]$.
            \item $j=5: A[5]=5 \le 8 \implies m=5$, troca $A[5] \leftrightarrow A[5]$.
            \item $j=6: A[6]=3 \le 8 \implies m=6$, troca $A[6] \leftrightarrow A[6]$.
            \item $j=7: A[7]=2 \le 8 \implies m=7$, troca $A[7] \leftrightarrow A[7]$.
            \item $j=8: A[8]=10 > 8 \implies$ nenhuma troca.
            \item $j=9: A[9]=9 > 8 \implies$ nenhuma troca.
        \end{itemize}
        \item Troca final: $A[p]=A[2]$ troca com $A[m]=A[7]$ (troca $8$ com $2$).
        \item Estado do vetor: $\langle 1, \mathbf{2, 7, 6, 5, 3}, \underline{8}, \mathbf{10, 9} \rangle$. Posição final: $q = 7$.
        \item Posto relativo: $k = 7 - 2 + 1 = 6$.
        \item Como $i = 2 < k = 6$: o elemento procurado está à esquerda. Ordem mantida: $i = 2$.
    \end{itemize}

    \item \textbf{Chamada 3:} \texttt{SELECTION(A, 2, 6, 2)}
    \begin{itemize}
        \item Subvetor: $A[2 \dots 6] = \langle \mathbf{2}, 7, 6, 5, 3 \rangle$, buscando $i = 2$.
        \item Pivô: $x = A[2] = 2$. Índice $m = 2$.
        \item Varredura ($j = 3 \dots 6$): todos os elementos $\{7, 6, 5, 3\}$ são maiores que $2$. Logo $m = 2$.
        \item Troca final: $A[2] \leftrightarrow A[2]$. Posição final: $q = 2$.
        \item Estado do vetor: $\langle 1, \underline{2}, \mathbf{7, 6, 5, 3}, 8, 10, 9 \rangle$.
        \item Posto relativo: $k = 2 - 2 + 1 = 1$.
        \item Como $i = 2 > k = 1$: o elemento está à direita. Nova ordem: $i' = 2 - 1 = 1$.
    \end{itemize}

    \item \textbf{Chamada 4:} \texttt{SELECTION(A, 3, 6, 1)}
    \begin{itemize}
        \item Subvetor: $A[3 \dots 6] = \langle \mathbf{7}, 6, 5, 3 \rangle$, buscando $i = 1$.
        \item Pivô: $x = A[3] = 7$. Índice $m$ inicializa em $3$.
        \item Varredura ($j = 4 \dots 6$):
        \begin{itemize}
            \item $j=4: A[4]=6 \le 7 \implies m=4$, troca $A[4] \leftrightarrow A[4]$.
            \item $j=5: A[5]=5 \le 7 \implies m=5$, troca $A[5] \leftrightarrow A[5]$.
            \item $j=6: A[6]=3 \le 7 \implies m=6$, troca $A[6] \leftrightarrow A[6]$.
        \end{itemize}
        \item Troca final: $A[p]=A[3]$ troca com $A[m]=A[6]$ (troca $7$ com $3$).
        \item Estado do vetor: $\langle 1, 2, \mathbf{3, 6, 5}, \underline{7}, 8, 10, 9 \rangle$. Posição final: $q = 6$.
        \item Posto relativo: $k = 6 - 3 + 1 = 4$.
        \item Como $i = 1 < k = 4$: o elemento está à esquerda. Ordem mantida: $i = 1$.
    \end{itemize}

    \item \textbf{Chamada 5:} \texttt{SELECTION(A, 3, 5, 1)}
    \begin{itemize}
        \item Subvetor: $A[3 \dots 5] = \langle \mathbf{3}, 6, 5 \rangle$, buscando $i = 1$.
        \item Pivô: $x = A[3] = 3$. Índice $m = 3$.
        \item Varredura ($j = 4 \dots 5$): elementos $\{6, 5\}$ são maiores que $3$. Nenhuma troca; $m = 3$.
        \item Troca final: $A[3] \leftrightarrow A[3]$. Posição final: $q = 3$.
        \item Estado do vetor: $\langle 1, 2, \underline{3}, 6, 5, 7, 8, 10, 9 \rangle$.
        \item Posto relativo: $k = 3 - 3 + 1 = 1$.
        \item Como $i = 1 == k = 1$: \textbf{condição de parada satisfeita!}
        \item O algoritmo retorna o elemento no pivô: $A[q] = A[3] = 3$.
    \end{itemize}
\end{enumerate}

\topico{Quadro Síntese das Chamadas}
\begin{center}
\small
\setlength{\tabcolsep}{4.5pt}
\begin{tabular}{ccccccc}
\toprule
\textbf{Chamada} & \textbf{Subvetor} & \textbf{Pivô ($x$)} & \textbf{Posição ($q$)} & \textbf{Posto ($k$)} & \textbf{Ordem ($i$)} & \textbf{Ação} \\ \midrule
1 & $A[1 \dots 9]$ & $A[1]=1$ & 1 & 1 & 3 & Recurso à direita ($i \leftarrow 2$) \\
2 & $A[2 \dots 9]$ & $A[2]=8$ & 7 & 6 & 2 & Recurso à esquerda ($i \leftarrow 2$) \\
3 & $A[2 \dots 6]$ & $A[2]=2$ & 2 & 1 & 2 & Recurso à direita ($i \leftarrow 1$) \\
4 & $A[3 \dots 6]$ & $A[3]=7$ & 6 & 4 & 1 & Recurso à esquerda ($i \leftarrow 1$) \\
5 & $A[3 \dots 5]$ & $A[3]=3$ & 3 & 1 & 1 & \textbf{Retorna $A[3] = 3$} \\
\bottomrule
\end{tabular}
\end{center}

\res{\text{O } 3^\circ \text{ menor elemento retornado pelo algoritmo é } 3.}

\vspace{12pt}
\hrule
\vspace{6pt}
% ==============================================================================
% QUESTÃO 3
% ==============================================================================
\questao{3}{Comparação de Métodos para Listar os $i$ Maiores Elementos em Ordem Crescente}

\textbf{Enunciado:} Dado um vetor $A[1, n]$, queremos imprimir em ordem crescente os $i$ maiores elementos de $A$. Compare a complexidade dos seguinte métodos em função de $n$ e $i$.
\begin{itemize}
    \item[(a)] Ordene o vetor e liste os $i$ maiores elementos.
    \item[(b)] Construa uma fila de prioridade (max-heap) e chame a rotina para extrair o máximo $i$ vezes.
    \item[(c)] Use um algoritmo de seleção para encontrar o $i$-ésimo maior elemento, particione o vetor em torno dele e ordene os $i$ maiores elementos.
\end{itemize}

\topico{Análise do Método (a): Ordenação Total}
\begin{enumerate}
    \item \textbf{Ordenação de todo o vetor $A[1 \dots n]$:} Utilizando um algoritmo ótimo baseado em comparações (como \texttt{MergeSort} ou \texttt{HeapSort}), o custo no pior caso é $\Theta(n \lg n)$.
    \item \textbf{Listagem dos $i$ maiores elementos:} Com o vetor ordenado de forma não decrescente, os $i$ maiores elementos ocupam as posições finais $A[n - i + 1 \dots n]$. Percorrer e imprimir essas posições consome tempo $\Theta(i)$.
\end{enumerate}
\[
T_a(n, i) = \Theta(n \lg n) + \Theta(i) = \Theta(n \lg n), \quad \text{pois } 1 \le i \le n.
\]

\topico{Análise do Método (b): Max-Heap e $i$ Extrações}
\begin{enumerate}
    \item \textbf{Construção da fila de prioridade (\texttt{BUILD-MAX-HEAP}):} Constrói o max-heap em tempo linear $\Theta(n)$.
    \item \textbf{Extração do máximo $i$ vezes:} Cada operação \texttt{EXTRACT-MAX} remove a raiz e restaura a propriedade de heap via \texttt{MAX-HEAPIFY} em tempo $O(\lg n)$ no pior caso. O custo acumulado das $i$ extrações é:
    \[
    \sum_{j=0}^{i-1} O(\lg(n - j)) = \Theta(i \lg n).
    \]
    \item \textbf{Ajuste para ordem crescente:} As extrações fornecem os elementos em ordem \textit{decrescente}. Armazená-los em um vetor auxiliar de trás para frente (ou inverter o resultado) custa $\Theta(i)$.
\end{enumerate}
\[
T_b(n, i) = \Theta(n) + \Theta(i \lg n) + \Theta(i) = \Theta(n + i \lg n).
\]

\topico{Análise do Método (c): Seleção Linear, Partição e Ordenação Parcial}
\begin{enumerate}
    \item \textbf{Seleção da estatística de ordem:} O $i$-ésimo maior elemento corresponde exatamente ao $(n - i + 1)$-ésimo menor elemento de $A$. Utilizando o algoritmo de seleção determinístico em tempo linear no pior caso (\texttt{SELECT} das medianas das medianas, Seção 9.3 do Cormen), esta etapa leva tempo $\Theta(n)$.
    \item \textbf{Particionamento:} Particionar o vetor em torno do elemento encontrado coloca os $i$ maiores elementos no subvetor final $A[n - i + 1 \dots n]$, com custo $\Theta(n)$.
    \item \textbf{Ordenação dos $i$ maiores elementos:} Ordenamos apenas o subvetor de tamanho $i$ em ordem crescente usando um algoritmo ótimo (\texttt{MergeSort}), com custo $\Theta(i \lg i)$.
    \item \textbf{Listagem:} Percorrer o subvetor já ordenado consome $\Theta(i)$.
\end{enumerate}
\[
T_c(n, i) = \Theta(n) + \Theta(n) + \Theta(i \lg i) + \Theta(i) = \Theta(n + i \lg i).
\]

\topico{Quadro Comparativo e Conclusão}
\begin{center}
\small
\setlength{\tabcolsep}{5pt}
\begin{tabular}{p{1.2cm}p{8.8cm}p{3.8cm}}
\toprule
\textbf{Método} & \textbf{Etapas Fundamentais} & \textbf{Complexidade} \\ \midrule
\textbf{(a)} & Ordenação total de $A$ + listagem dos $i$ maiores & $\Theta(n \lg n)$ \\
\textbf{(b)} & \texttt{BUILD-MAX-HEAP} + $i \times$ \texttt{EXTRACT-MAX} + reversão & $\Theta(n + i \lg n)$ \\
\textbf{(c)} & \texttt{SELECT} linear + partição + ordenação de $i$ elementos & $\Theta(n + i \lg i)$ \\
\bottomrule
\end{tabular}
\end{center}

Como $1 \le i \le n$, vale sempre $\lg i \le \lg n$, de modo que:
\[
\Theta(n + i \lg i) \le \Theta(n + i \lg n) \le \Theta(n \lg n).
\]
\begin{itemize}
    \item Para $i = O(1)$ ou $i = O(n/\lg n)$: os métodos (b) e (c) operam em tempo linear $\Theta(n)$, superando assintoticamente o método (a).
    \item Para $i = \Theta(n)$: todos os três métodos têm ordem assintótica $\Theta(n \lg n)$, mas o método (c) apresenta constante proporcional menor ($\alpha \lg(\alpha n) < \lg n$).
\end{itemize}
\res{\text{O Método (c) é assintoticamente o mais eficiente: } T_c(n, i) = \Theta(n + i \lg i).}

\vspace{12pt}
\hrule
\vspace{6pt}
% ==============================================================================
% QUESTÃO 4
% ==============================================================================
\questao{4}{QuickSort Modificado com Tempo $\Theta(n \lg n)$ no Pior Caso}

\textbf{Enunciado:} Mostre como modificar \texttt{QuickSort} de modo que tenha complexidade de tempo $\Theta(n \lg n)$ no pior caso.

\topico{Diagnóstico do QuickSort Convencional}
O \texttt{QuickSort} clássico exibe tempo $\Theta(n^2)$ no pior caso quando as partições são severamente desbalanceadas (por exemplo, subproblemas de tamanhos $0$ e $n-1$), resultando na relação $T(n) = T(n-1) + \Theta(n) = \Theta(n^2)$.

\topico{A Modificação Proposta}
Para eliminar o desbalanceamento no pior caso, selecionamos a \textbf{mediana exata} do subvetor como pivô antes de cada particionamento. 
Como estabelecido pelo algoritmo \texttt{SELECT} (Cormen et al., Seção 9.3 -- algoritmo da mediana das medianas), a mediana de qualquer conjunto de $m$ elementos pode ser encontrada no pior caso em tempo linear $O(m)$.

\topico{Pseudocódigo do QuickSort Modificado}
\begin{lstlisting}
QUICKSORT-PIOR-CASO-OTIMO(A, p, r)
 1  se p < r então
 2      m <- r - p + 1
 3      k <- piso((m + 1) / 2)               // posto da mediana
 4      med <- SELECT(A, p, r, k)            // determina a mediana em tempo O(m)
 5      localiza indice idx em [p, r] onde A[idx] == med
 6      troca A[idx] com A[r]               // posiciona a mediana como pivo
 7      q <- PARTITION(A, p, r)              // particiona em torno da mediana
 8      QUICKSORT-PIOR-CASO-OTIMO(A, p, q - 1)
 9      QUICKSORT-PIOR-CASO-OTIMO(A, q + 1, r)
\end{lstlisting}

\topico{Demonstração da Complexidade $\Theta(n \lg n)$ no Pior Caso}
Ao utilizar a mediana exata como pivô em um subvetor de tamanho $m$:
\begin{itemize}
    \item O subvetor à esquerda tem tamanho no máximo $\lfloor (m - 1)/2 \rfloor \le m/2$.
    \item O subvetor à direita tem tamanho no máximo $\lceil (m - 1)/2 \rceil \le m/2$.
\end{itemize}
Portanto, a partição divide o problema em dois subproblemas de tamanho no máximo $n/2$. 
O trabalho não recursivo em cada nível é composto por:
\begin{itemize}
    \item Execução de \texttt{SELECT}: $O(n)$ no pior caso;
    \item Localização do índice e troca: $O(n)$;
    \item Execução de \texttt{PARTITION}: $\Theta(n)$.
\end{itemize}
Logo, o custo não recursivo total é limitado por $c \cdot n$, para alguma constante $c > 0$. A recorrência do pior caso satisfaz:
\[
T(n) \le 2\,T\left(\frac{n}{2}\right) + c\,n.
\]
Aplicando o \textbf{Teorema Mestre}:
\begin{itemize}
    \item Parâmetros: $a = 2$, $b = 2$, $f(n) = c n$.
    \item Expoente crítico: $n^{\log_b a} = n^{\log_2 2} = n^1 = n$.
    \item Como $f(n) \in \Theta(n^{\log_b a}) = \Theta(n)$, estamos no \textbf{Caso 2} do Teorema Mestre.
\end{itemize}
Dessa forma:
\[
T(n) \in \Theta(n^1 \lg n) = \Theta(n \lg n).
\]
Como qualquer algoritmo de ordenação por comparações possui cota inferior $\Omega(n \lg n)$ no pior caso, concluímos:
\res{T(n) \in \Theta(n \lg n) \quad \text{no pior caso.}}

\vspace{12pt}
\hrule
\vspace{6pt}
% ==============================================================================
% QUESTÃO 5
% ==============================================================================
\questao{5}{Algoritmo Linear de Seleção Usando Caixa-Preta de Mediana}

\textbf{Enunciado:} Suponha que você tenha uma subrotina do tipo ``caixa-preta'' que determina a mediana em tempo linear (no pior caso). Descreva um algoritmo linear simples que resolve o problema da seleção para todo $i$.

\topico{Descrição do Algoritmo}
Seja \texttt{MEDIANA(A, p, r)} a sub-rotina caixa-preta que retorna a mediana do subvetor $A[p \dots r]$ em tempo $c_1 \cdot m$, onde $m = r - p + 1$. O algoritmo \texttt{SELECAO-LINEAR} utiliza essa mediana como pivô perfeito para o particionamento, descartando garantidamente metade dos elementos a cada passo recursivo.

\topico{Pseudocódigo}
\begin{lstlisting}
SELECAO-LINEAR(A, p, r, i)
 1  se p == r então
 2      retorne A[p]
 3  med <- MEDIANA(A, p, r)                 // caixa-preta: custo O(m) no pior caso
 4  localiza indice idx em [p, r] tal que A[idx] == med
 5  troca A[idx] com A[r]                   // posiciona o pivo no final
 6  q <- PARTITION(A, p, r)                 // particiona em torno da mediana
 7  k <- q - p + 1                          // posto relativo do pivo
 8  se i == k então
 9      retorne A[q]
10  senão se i < k então
11      retorne SELECAO-LINEAR(A, p, q - 1, i)
12  senão
13      retorne SELECAO-LINEAR(A, q + 1, r, i - k)
\end{lstlisting}

\topico{Análise da Complexidade de Tempo no Pior Caso}
Como o pivô escolhido é a mediana exata de $A[p \dots r]$, o número de elementos estritamente menores que a mediana é no máximo $\lfloor (m - 1)/2 \rfloor \le m/2$, e o número de elementos estritamente maiores é no máximo $\lceil (m - 1)/2 \rceil \le m/2$.

Diferentemente do \texttt{QuickSort}, o algoritmo de seleção realiza \textbf{apenas uma chamada recursiva} (ou no lado esquerdo ou no lado direito), reduzindo o tamanho da entrada para no máximo $\lfloor n/2 \rfloor \le n/2$.

O custo não recursivo de cada etapa é dominado por:
\begin{itemize}
    \item Chamada a \texttt{MEDIANA}: $c_1 n$;
    \item Busca linear de \texttt{idx} e troca: $c_2 n$;
    \item Particionamento linear \texttt{PARTITION}: $c_3 n$.
\end{itemize}
Somando essas parcelas, o trabalho por etapa é limitado por $d \cdot n$ para uma constante $d > 0$. Assim, a relação de recorrência para o pior caso é:
\[
T(n) \le T\left(\frac{n}{2}\right) + d\,n.
\]

\topico{Resolução da Recorrência por Expansão (Série Geométrica)}
Expandindo a relação sucessivamente:
\begin{align*}
T(n) &\le d n + T(n/2) \\
     &\le d n + d\left(\frac{n}{2}\right) + T(n/4) \\
     &\le d n \left(1 + \frac{1}{2} + \frac{1}{4} + \dots + \frac{1}{2^{\lfloor \lg n \rfloor}}\right) + \Theta(1).
\end{align*}
Limitando a soma finita pela série geométrica infinita convergente de razão $1/2$:
\[
T(n) \le d n \sum_{j=0}^{\infty} \left(\frac{1}{2}\right)^j = d n \cdot \left(\frac{1}{1 - 1/2}\right) = 2 d n.
\]
Portanto, $T(n) \le 2 d n = O(n)$. 

Como qualquer algoritmo que resolve o problema da seleção precisa examinar os $n$ elementos no pior caso ($T(n) \in \Omega(n)$), temos:
\res{T(n) \in \Theta(n) \quad \text{no pior caso, para todo } i.}

\end{document}
